\documentclass[pre,twocolumn,superscriptaddress,nofootinbib]{revtex4-2}
\usepackage{graphicx}
\usepackage{amsmath}
\usepackage{booktabs}
\usepackage[hidelinks]{hyperref}

\begin{document}

\title{Does concentrated energy re-curl space? Three sealed tests in a graph model of emergent geometry, and a melt
every time}

\author{Emily Smith}
\email{emilysmithcreate@gmail.com}
\affiliation{Independent researcher, Charlottesville, Virginia, USA}

\date{DRAFT, 25 September 2026}

\begin{abstract}
In the family of graph energies obtained from combinatorial quantum gravity (CQG) by raising the coefficient
$\lambda$ of its local term above 1 (only $\lambda = 1$ is CQG; everything here is at $\lambda = 1.25$, next to CQG
and not in it), the flat torus is the exact ground state and a curled direction costs $4(\lambda - 1)$ per vertex.
The author's hypothesis holds that a black hole is a region of space re-curled toward the arrangement it came from,
and that concentrated energy is what pays for the re-curling. We ask the model, sealed and energy-conserving, what
energy does to flat space, three ways: spread through a bath any move may draw on, packed into one place as heat only
one vertex may spend or as disorder wired into one patch, and spread then drained at four rates from a slow leak to
a quench. Read by one rule fixed before the runs, vertices whose local dimension falls below 2 are folded and those
whose dimension rises above 2 are melted. In every one of the 22 cells of the local test, the 15 of the leaking test,
and the earlier spread test, at 36 to 256 vertices, with named points and with the points treated as
interchangeable, the majority outcome is a melt; no cell folds; the largest folded piece anywhere is eleven vertices;
and a melt neither folds nor flattens as its energy drains. Two dimensions do not re-curl under this energy. Every
verdict quoted was pre-registered before its run.
\end{abstract}

\maketitle

\section{Introduction}

The companion papers \cite{paper1,paper2} study a curled torus opening into flat space and what it leaves behind. This
one asks the reverse: can flat space be curled back up? The author's framework needs it to be possible, because in
her picture a black hole is a region re-curled toward the curled arrangement, and the energy concentrated in it is
what pays the curling cost \cite{repo}. The published programme reads a black hole the other way, as a bubble that has
reverted to the random phase \cite{T25}. The two readings differ in what concentrated energy produces, a fold or a
melt, and that is measurable.

Two facts fix the terms. Above $\lambda = 1$ the flat torus is the exact ground state, nothing lies below it, and
leaving it costs at least 32 by any single move (Sec.~II). And the two kinds of damage have different prices per
square: at $\lambda = 1.25$ a third square added to an edge costs 1, a square destroyed costs 16, so curling is the
cheaper damage by a factor of sixteen per square. A budget that cannot afford to melt can still afford to curl. The
question is which the dynamics buys.

\section{Model, the sealed box, and the reading}

The state is a four-regular bipartite graph on $N$ vertices under the hard-core rule of CQG. With $S$ squares and
$X = \sum_e (S_e - 2)_+$ surplus squares,
\begin{equation}
  H = 16(N - S) + 4\lambda X ,
\end{equation}
so a square destroyed costs 16 and a surplus square costs $4\lambda$. The local dimension $d(v)$ is the number of the
six edge pairs at $v$ that close no square: 2 on the flat sheet, 1 on a curled direction, 0 in a 4-cube, above 2 where
squares are missing. A vertex is \emph{folded} if $d < 2$ and \emph{melted} if $d > 2$; a replica's outcome, read from
its final block, is FOLDED if at least four vertices are damaged and folded ones are at least as many as melted,
MELTED if at least four are damaged and melted ones outnumber folded, HEALED if fewer than four are damaged. Four is
the smallest curled object, one column. A cell's outcome is the majority over replicas.

Sealed dynamics keep stores of energy beside the graph \cite{Creutz83}: a proposed switch draws on one store, is
refused if the store cannot pay, and pays into it when it lowers the energy; $H$ plus the stores is conserved to the
last unit in every run reported here. Three protocols differ only in which store a move may use and where the energy
starts. \emph{Spread}: a shared bath of $2N$ stores, the budget in one of them, any move anywhere may draw on it
(the earlier test, T21, and its single-demon predecessor). \emph{Packed}: either the sheet perfect and the budget in
the store of one vertex under a per-vertex bath, in which a move pays only from the store of the first vertex it
involves; or the budget wired into the graph as disorder within a radius of one vertex, the bath empty. \emph{Drained}:
the spread protocol with every store scaled by $1 - \ell$ after each block of 100 sweeps, at $\ell = 0.001$, 0.01,
0.1 and 1.0, the last a quench. Under the per-vertex bath a store below 32 can never act, because leaving the flat
sheet costs at least that, so the packed budgets are 32 to 192 units in one store, or 288 and 576 spread over nine.

With the points treated as interchangeable, every arrangement is weighted by its number of symmetries (the
relabelings that keep every edge), which favors a complete fold, a 4-cube with 192 symmetries, over any partial
state. The packed test was run that way at 64 vertices, the smallest size at which a complete fold exists, with the
fast automorphism count of Ref.~\cite{repo}.

\section{Spread through a bath: it melts}

\emph{T21, pre-registered.} At $N = 36$ and 64, budgets of 2, 4, 8 and 16 per vertex through the bath and 0.5 to 4
per vertex through a single demon, six replicas each, named and interchangeable points side by side. Every cell at
64 vertices has a melted majority under both weightings; at 36 vertices under the bath the interchangeable points
lean toward folding in one cell and reverse in the next, which the pre-registered rule reads as fragile. The verdict
was MELTS EITHER WAY at $N = 64$. This was the result that prompted the author's objection: a black hole is not
energy spread evenly but energy packed into one place, and heating a whole sheet at equilibrium randomizes it by
construction.

\section{Packed into one place: it melts}

\emph{T26, pre-registered, with the author's prediction that it re-curls.} Twenty-two cells, twenty replicas each at
144 and 256 vertices and ten at 64, 50,000 sweeps (4,000 at 64), the final graph saved. In the per-vertex protocol the
first move a store can buy is the cheapest exit from the sheet, a loss of two squares, which leaves eight vertices at
$d = 3$: melted, by the reading. What follows is a local heat cycle, the defect healing back into the store and being
bought again, and it never turns into a fold. Table~\ref{tab:t26} gives the census. No cell has a folded majority;
the largest folded piece in any final state is eleven vertices, at a budget of 96 in one store; the share of replicas
that melt first and fold later is at most one in twenty. A hot patch in a cold shared bath mostly stays melted at 144
and 256 vertices, healing in 5 to 20\% of replicas; at 64 vertices the named patch melts in 7 of 10 and the
interchangeable one heals in 10 of 10 at 32 units and melts in 7 of 10 at 64. The verdict is MELTS. The author's
prediction, confirmed as hers before any cell was read, failed.

\begin{table}[b]
\caption{The packed test (T26): outcomes over replicas, from the final block. Stores: the budget in one vertex's store
(radius 0) or nine (radius 2). Patch: disorder wired within radius 3 (2 at $N = 64$), cold shared bath; ``local heat''
keeps the heat the patch gives off in per-vertex stores.}
\label{tab:t26}
\begin{ruledtabular}
\begin{tabular}{llrrrr}
protocol & $N$ & budgets & melted & folded & healed \\
\hline
stores, radius 0 & 144 & 32 to 192 & 115 & 4 & 1 \\
stores, radius 0 & 256 & 32 to 192 & 119 & 0 & 1 \\
stores, radius 2 & 144 & 288, 576 & 40 & 0 & 0 \\
patch, shared bath & 144 & 32, 64, 128 & 48 & 3 & 9 \\
patch, shared bath & 256 & 32, 64, 128 & 57 & 0 & 3 \\
patch, local heat & 144 & 32, 64, 128 & 60 & 0 & 0 \\
patch, named & 64 & 32, 64 & 14 & 4 & 2 \\
patch, interchangeable & 64 & 32, 64 & 7 & 0 & 13 \\
\end{tabular}
\end{ruledtabular}
\end{table}

\section{Drained: it stays melted}

\emph{T27, pre-registered, with the author's prediction that the melt folds before it flattens.} The spread protocol
at 144 vertices, budgets of 2, 4 and 8 per vertex, drained at four rates and a sealed control, twelve replicas each,
30,000 sweeps. All fifteen cells have a melted majority: 57 of the 60 draining replicas melted, the largest folded
piece anywhere 16 vertices. The quench freezes the disorder as it stands; the slow drains do not anneal it within the
run. The verdict is STAYS MELTED; ours held, the author's failed.

\section{Discussion}

Three protocols, two weightings of the points, sizes from 36 to 256 vertices: under this energy and this dynamics,
energy given to flat two-dimensional space melts it and does not curl it, however it is delivered and however it is
taken away. Curling is the cheaper damage per square, and the dynamics does not buy it, because the moves that lower
squares are the ones the chain offers first from a flat sheet, and no single move adds a third square to an edge of
the perfect sheet (Sec.~II of Ref.~\cite{paper1}). A fold would have to be reached through melt, and the melt does
not go on to fold.

What this does and does not show. It does not show that the author's black hole is impossible: her rule is that the
three space directions curl together, so the fold she means is three-dimensional, and this paper is two-dimensional.
It shows that the two-dimensional version of the question, asked three ways, has one answer, and that a re-curling
mechanism, if it exists in this family, needs either three directions or an ingredient this energy lacks. The
six-link version of the packed test is the next thing to run, once the six-link runs have shown what a curled
region looks like there \cite{repo}. And it carries the model's standing caveat: the melt here is what Ref.~\cite{T25}
calls matter, and the author does not; what the melt is called does not change what it does.

\begin{acknowledgments}
Code, analysis and text were drafted with Claude (Anthropic) in conversation with the author, who directed and
checked the work; no physicist has reviewed it. The repository holds every configuration, seed, result and
pre-registration: \url{https://github.com/EmilySmithCreate/RecreationalPhysics}.
\end{acknowledgments}

\begin{thebibliography}{9}
\bibitem{paper1} E. Smith, The curled torus burps: front-driven decompactification between ordered states in a
graph model of emergent geometry, draft (2026), \texttt{docs/papers/curled\_torus/}.
\bibitem{paper2} E. Smith, What the burp leaves behind, draft (2026), \texttt{docs/papers/relic/}.
\bibitem{T25} C. A. Trugenberger, Combinatorial quantum gravity and emergent 3D quantum behaviour, review (2025);
arXiv:2512.17676.
\bibitem{Creutz83} M. Creutz, Microcanonical Monte Carlo simulation, Phys. Rev. Lett. 50, 1411 (1983).
\bibitem{repo} \url{https://github.com/EmilySmithCreate/RecreationalPhysics}: PREREGISTRATION T21, T26, T27;
ASSUMPTIONS O20, O22, O36, O48, Q20, Q22.
\end{thebibliography}

\end{document}
